% Generated draft: requires mathematical and source review.
\documentclass[11pt]{article}
\usepackage[margin=25mm]{geometry}
\usepackage{fontspec}
\setmainfont{Latin Modern Roman}
\usepackage{amsmath,amssymb,mathtools,unicode-math}
\setmathfont{Latin Modern Math}
\usepackage{enumitem,needspace,xurl,hyperref}
\setlength{\parindent}{0pt}
\setlength{\parskip}{5pt}
\setlength{\emergencystretch}{3em}
\newcommand{\N}{\mathbb{N}}
\newcommand{\Z}{\mathbb{Z}}
\newcommand{\Q}{\mathbb{Q}}
\newcommand{\R}{\mathbb{R}}
\newcommand{\C}{\mathbb{C}}
\newcommand{\F}{\mathbb{F}}
\newcommand{\A}{\mathbb{A}}
\newcommand{\PP}{\mathbb{P}}
\newcommand{\E}{\mathbb{E}}
\newcommand{\eps}{\varepsilon}
\newcommand{\dd}{\,\mathrm{d}}
\begin{document}
\section*{UnsolvedMath ID: 701207}
\textbf{Source problem code:} catalogue entry
\subsection{Definitions and mathematical statement}
Let \(G\) be a group with identity \(1\). Use the commutator convention
\[
[x,y]=x^{-1}y^{-1}xy,\qquad [h,{}_1g]=[h,g],\qquad [h,{}_{n+1}g]=[[h,{}_ng],g]\quad(n\geq1).
\]
An element \(g\in G\) is an \emph{Engel element} if
\[
\forall h\in G\ \exists k\in\Z_{\geq1}:\ [h,{}_kg]=1.
\]
It is \emph{boundedly Engel} if
\[
\exists k\in\Z_{\geq1}\ \forall h\in G:\ [h,{}_kg]=1.
\]
The bound may depend on \(g\) and \(G\). Write
\[
B(G)=\{g\in G:g\text{ is boundedly Engel}\}.
\]
A group is \emph{torsion-free} if \(x^n=1\), for \(x\in G\) and \(n\in\Z_{\geq1}\), implies \(x=1\). A group is \emph{right-ordered} if it admits a strict total order \(<\) such that \(x<y\) implies \(xz<yz\) for all \(x,y,z\in G\). Here \emph{linearly ordered group} is interpreted in the usual group-theoretic sense of a group admitting a strict total order invariant under both left and right multiplication: \(x<y\) implies \(axb<ayb\) for all \(a,b,x,y\in G\). The supplied source does not explicitly specify its order or commutator conventions; these conventions make the draft precise and require confirmation against the source.

The remaining questions of Kourovka Notebook Problem 16.15 are:
\begin{enumerate}
\item[(b)] Is \(B(G)\) a subgroup of \(G\) for every torsion-free group \(G\)?
\item[(c)] Is \(B(G)\) a subgroup of \(G\) for every right-ordered group \(G\)?
\item[(d)] Is \(B(G)\) a subgroup of \(G\) for every linearly ordered group \(G\), with the interpretation specified above?
\end{enumerate}
The archived part (a), attributed specifically to B. I. Plotkin, asks whether \(B(G)\) is a subgroup for every group \(G\), without additional hypotheses. The supplied archive records a negative answer to (a), attributed to A. I. Sozutov. This is a supplied archival assessment, not an independently verified result here. The registry labels the overall problem partially solved and retains (b)--(d) as the remaining questions.
\subsection{Short English statement}
For each of three classes of groups---torsion-free groups, right-ordered groups, and groups with an order invariant under multiplication on both sides---ask whether the boundedly Engel elements form a subgroup. Each element may have its own Engel bound, but that bound must work for every other element of the group. The supplied archive reports that the unrestricted question (a) has a negative answer; the registry retains (b)--(d) as unresolved.
\subsection{Sources}
\begin{itemize}
\item Local UnsolvedMath registry, supplied complete record with permanent ID \(701207\), exact source code \texttt{catalogue entry}, title \emph{Kourovka Notebook Problem 16.15}; upstream revision \texttt{d371ee917ac9647fbb06cd6fc529deec6f3dda82}.
\item \emph{The Kourovka Notebook (Unsolved Problems in Group Theory)}, 16th issue (2006), Problem 16.15; proposed by V. V. Bludov. The supplied background identifies an update in the 21st issue (2026), \texttt{arXiv:1401.0300v48}. The supplied collection description names E. I. Khukhro and V. D. Mazurov as editors.
\item Archived part (a) is attributed in the supplied statement to B. I. Plotkin. The supplied Archive of Solved Problems entry reports: negative answer, A. I. Sozutov, \emph{Siberian Math. J.}, \textbf{60}, no. 6 (2019), 1099--1100. No article title or DOI is supplied.
\item The supplied status review is dated \texttt{2026-10-16}, labels the record \texttt{partially\_solved}, and states that part (a) is solved in the archive while the remaining parts occur in the main section. This is supplied dated curation, not an independently verified literature assessment.
\end{itemize}
\end{document}
